\section{Fourier concentration of squares under product averages}
\label{sec:fourier}

We next show how averaging a positive semidefinite matrix-valued
function can produce a representation by squares whose Fourier
coefficients concentrate around a separate center for each square.
The assumption is a scalar variance bound for each coordinate average.
This bound extends to matrix-valued functions with the same constant.
We will use it to control a fourth-order Gram potential and then
convert that bound into an estimate on the Fourier coefficients.

Fix an integer $a\ge1$, put $L=2^a$, and identify each of
$Y_1,\ldots,Y_t$ with the additive group $(\mathbb F_2)^a$.
We use a second copy $\Lambda=(\mathbb F_2)^a$ to index its characters.
For $p=(p_1,\ldots,p_t)\in\Lambda^t$ and
$y=(y_1,\ldots,y_t)\in\prod_iY_i$, define
\[
  \chi_p(y)=(-1)^{\sum_{i=1}^t p_i\cdot y_i},
\]
where the inner products and the exponent are computed modulo two.
These real characters form an orthonormal basis for uniform probability
on $\prod_iY_i$.  For a real matrix-valued function $H$, write
\[
  \widehat H_p=\E_y[\chi_p(y)H(y)],
  \qquad H(y)=\sum_{p\in\Lambda^t}\widehat H_p\chi_p(y).
\]
The block degree of $\chi_p$ is $|\{i:p_i\ne0\}|$.
A function has block degree at most $D$ if its Fourier coefficients
vanish at every character of block degree greater than $D$; this
convention includes the zero function.
We also write
$\operatorname{dist}(p,q)=|\{i:p_i\ne q_i\}|$.
Multiplication by $\chi_q$ translates frequency $p$ to $p+q$, whose
block degree is $\operatorname{dist}(p,q)$.  In particular, a real
function supported within distance $D$ of $q$ becomes a function of
block degree at most $D$ after this multiplication, with its square
unchanged.  This translation of Fourier support while preserving
squares also appears in the construction of sparse sums of squares
on finite abelian groups by Fawzi, Saunderson, and
Parrilo~\cite{fsp-sparse}.

\begin{theorem}[Fourier smoothing of matrix squares]
\label{thm:fourier-smoothing}
Let $t,r\ge1$ be integers.  For each $1\le i\le t$, let $X_i$ and
$\mathcal M_i$ be nonempty finite sets with uniform probability, and
let $Y_i$ be a specified copy of $(\mathbb F_2)^a$, where $a\ge1$ and
$L=2^a$.  For $m\in\mathcal M_i$ and $y\in Y_i$, let
$\pi_{i,y}^m$ be a probability distribution on $X_i$.
Suppose that, for some $\theta\ge0$, these distributions satisfy
\begin{equation}\label{eq:channel-marginal}
  \E_{m\in\mathcal M_i}\pi_{i,y}^m(x)=\frac1{|X_i|}
  \qquad(x\in X_i, y\in Y_i)
\end{equation}
and, for every real function $g:X_i\to\R$ and every $y\in Y_i$,
\begin{equation}\label{eq:channel-variance}
  \E_{m\in\mathcal M_i}
  \left(\sum_{x\in X_i}\pi_{i,y}^m(x)g(x)-\E_{x\in X_i}g(x)\right)^2
  \le\theta\E_{x\in X_i}g(x)^2.
\end{equation}
Let $F:\prod_iX_i\to\mathbb S_+^r$ satisfy $\tr F(x)\le B$ for all
$x$, where $B>0$ and $\mathbb S_+^r$ is the cone of real symmetric
positive semidefinite $r$-by-$r$ matrices.  Fix $\lambda>1$ such that
\begin{equation}\label{eq:smoothing-parameter}
  L\lambda^2C_L\theta\le1,
\end{equation}
where $C_L=2^{L-1}$ is the constant of Lemma~\ref{lem:splitting}.
For $m=(m_1,\ldots,m_t)\in\prod_i\mathcal M_i$, define
\begin{equation}\label{eq:product-average}
  \overline F_m(y)
  =\sum_{x\in\prod_iX_i}
     \left(\prod_{i=1}^t\pi_{i,y_i}^{m_i}(x_i)\right)F(x).
\end{equation}

There is a finite integer $R$ and a family of functions
$H_q^m:\prod_iY_i\to\R^{R\times r}$, indexed by
$m\in\prod_i\mathcal M_i$ and $q\in\Lambda^t$, such that
\begin{equation}\label{eq:smoothing-squares}
  \overline F_m(y)=\sum_{q\in\Lambda^t}H_q^m(y)^{\mathsf T}H_q^m(y).
\end{equation}
For each $m,q$, form the horizontal concatenation and its Gram matrix
\begin{equation}\label{eq:smoothing-gram}
  U_q^m=
    \left[\lambda^{\operatorname{dist}(q,p)/2}
           \widehat H_{q,p}^m\right]_{p\in\Lambda^t},
  \qquad C_q^m=U_q^m(U_q^m)^{\mathsf T},
\end{equation}
where $\widehat H_{q,p}^m$ is the Fourier coefficient of $H_q^m$ at $p$.
The family can be chosen so that, for independent uniform $m_i$,
\begin{equation}\label{eq:smoothing-potential}
  \E_m\sum_{q\in\Lambda^t}\normF{C_q^m}^2\le2^tB^2.
\end{equation}
Consequently, at least one $m$ satisfies
\begin{equation}\label{eq:smoothing-coefficients}
  \sum_{q,p\in\Lambda^t}
    \lambda^{\operatorname{dist}(q,p)}\normF{\widehat H_{q,p}^m}^2
  \le L^t\sqrt r\,B\,2^{t/2}.
\end{equation}
\end{theorem}

The weights in \eqref{eq:smoothing-coefficients} penalize frequencies
far from the center $q$ of their branch.  To explain the potential,
fix $m,q$, suppress these indices, and put
$w_p=\lambda^{\operatorname{dist}(q,p)}$.  Then
\[
  \tr C=\sum_p w_p\normF{\widehat H_p}^2,
  \qquad
  \normF{C}^2=\sum_{p,p'}w_pw_{p'}
          \normF{\widehat H_p^{\mathsf T}\widehat H_{p'}}^2.
\]
Thus the trace measures the weighted Fourier mass we want to bound,
whereas the squared Frobenius norm measures squared overlaps of
coefficient matrices.  The latter quantity matches the squared norm
in the variance assumption.

We prove the theorem by processing the coordinates in order.  A direct
sum of squares realizes each new average, and
Lemma~\ref{lem:splitting} separates its new frequencies into branches.
We bound the increase of the Gram potential at each step, then use
the ranks of the final Gram matrices to bound their traces.

\begin{proof}
For a matrix-valued function $Z$ on $X_i$, use the notation
\[
  T_{i,y}Z(m)=\sum_{x\in X_i}\pi_{i,y}^m(x)Z(x).
\]
Applying \eqref{eq:channel-variance} to each entry and summing gives
\begin{equation}\label{eq:matrix-channel-variance}
  \E_m\normF{T_{i,y}Z(m)-\E_xZ(x)}^2
  \le\theta\E_x\normF{Z(x)}^2.
\end{equation}
This inequality holds in every finite matrix dimension, with the same
$\theta$.

\paragraph{Partial averages and the potential.}
Use the abbreviations $m_{\le i}=(m_1,\ldots,m_i)$ and
$x_{>i}=(x_{i+1},\ldots,x_t)$, and likewise for $y$.
At stage $i$ we construct functions
\[
  H_q^{(i)}(y_{\le i};x_{>i}),\qquad q\in\Lambda^i,
\]
depending also on $m_{\le i}$, such that
\begin{equation}\label{eq:partial-squares}
  \sum_{q\in\Lambda^i}(H_q^{(i)})^{\mathsf T}H_q^{(i)}
  =\sum_{x_{\le i}}
       \left(\prod_{j=1}^i\pi_{j,y_j}^{m_j}(x_j)\right)F(x).
\end{equation}
All matrices at stage $i$ have a common row dimension $R_i$ and $r$
columns.  Their Fourier coefficients refer only to the processed
variables $y_{\le i}$.  Define
\begin{equation}\label{eq:partial-gram}
  U_q^{(i)}=
    \left[\lambda^{\operatorname{dist}(q,p)/2}
           \widehat H_{q,p}^{(i)}\right]_{p\in\Lambda^i},
  \qquad C_q^{(i)}=U_q^{(i)}(U_q^{(i)})^{\mathsf T}.
\end{equation}
Thus $U_q^{(i)}$ has $rL^i$ columns and is a function of $x_{>i}$.
The potential is
\[
  \Phi_i=\E_{m_{\le i},x_{>i}}
          \sum_{q\in\Lambda^i}\normF{C_q^{(i)}}^2,
\]
where all listed variables are independent and uniform.
At stage zero take $R_0=r$ and $H_{\varnothing}^{(0)}=F^{1/2}$.
Then $C_{\varnothing}^{(0)}=F$, so
$\Phi_0\le\E_x(\tr F(x))^2\le B^2$.
We will prove $\Phi_{i+1}\le2\Phi_i$ while preserving
\eqref{eq:partial-squares}.

\paragraph{Realizing one more average.}
Fix $m_{\le i}$, a branch $q\in\Lambda^i$, and $x_{>i+1}$.
For the new coordinate abbreviate $X_{i+1},Y_{i+1},m_{i+1}$ by
$X,Y,m$, and write its input as $x$ and its character variable as
$y$.  The old concatenation is now a matrix function
$U(x)=U_q^{(i)}(x;x_{>i+1})$.
Put
\[
  Z(x)=U(x)^{\mathsf T}U(x).
\]
It has size $rL^i$ by $rL^i$, and
\begin{equation}\label{eq:two-gram-norms}
  \normF{Z(x)}^2=\normF{C_q^{(i)}(x)}^2.
\end{equation}
The fixed data determine the entire function $x\mapsto U(x)$ before
the new averaging parameter $m$ is chosen.

For a fixed $m$, form a matrix-valued function $H'$ with
$R_{i+1}=L|X|R_i$ rows.  We give different values of $y$ disjoint row
blocks by indexing those blocks with $(y',x)\in Y\times X$.
This extra label will eliminate products between different $y$-values
from the Fourier overlap calculation below.
The block at $(y',x)$ is
\begin{equation}\label{eq:tagged-stack}
  H'(y_{\le i},y)_{(y',x)}
  =\mathbf1_{\{y'=y\}}\sqrt{\pi_y^m(x)}\,
       H_q^{(i)}(y_{\le i};x,x_{>i+1}).
\end{equation}
Consequently,
\[
  H'^{\mathsf T}H'
  =\sum_{x\in X}\pi_y^m(x)
       H_q^{(i)}(y_{\le i};x,x_{>i+1})^{\mathsf T}
       H_q^{(i)}(y_{\le i};x,x_{>i+1}).
\]
This stack performs the new average.  We next split its space of row
coordinates in a way that controls its Fourier coefficients.

For $z\in\Lambda$, concatenate the Fourier coefficients of $H'$ at
$(p,z)$, using only the old weights:
\[
  V_z=\left[\lambda^{\operatorname{dist}(q,p)/2}
                 \widehat H'_{p,z}\right]_{p\in\Lambda^i},
  \qquad A_z=V_zV_z^{\mathsf T}.
\]
The $(y',x)$ row block of $V_z$ is
$L^{-1}\chi_z(y')\sqrt{\pi_{y'}^m(x)}U(x)$.
Multiplying the stacked matrices therefore gives the identity
\begin{equation}\label{eq:new-frequency-overlap}
  V_z^{\mathsf T}V_{z'}
  =\frac1{L^2}\sum_{y\in Y}
       \chi_z(y)\chi_{z'}(y)T_{i+1,y}Z(m).
\end{equation}
In particular,
$\tr(A_zA_{z'})=\normF{V_z^{\mathsf T}V_{z'}}^2$.

\paragraph{The overlap estimate.}
For $z=z'$, Jensen's inequality and
\eqref{eq:channel-marginal} applied to
\eqref{eq:new-frequency-overlap} give
\[
  \E_m\tr(A_z^2)
  \le\frac1{L^2}\E_x\normF{Z(x)}^2.
\]
For $z\ne z'$, let $\mu=\E_xZ(x)$.
Character orthogonality permits subtracting $\mu$ from every summand
in \eqref{eq:new-frequency-overlap}.  Jensen's inequality and
\eqref{eq:matrix-channel-variance} now give
\[
\begin{split}
  \E_m\tr(A_zA_{z'})
  &\le\frac1{L^2}\E_y\E_m
           \normF{T_{i+1,y}Z(m)-\mu}^2\\
  &\le\frac\theta{L^2}\E_x\normF{Z(x)}^2.
\end{split}
\]
Together with \eqref{eq:two-gram-norms}, these estimates say that
\begin{equation}\label{eq:overlap-summary}
  \E_m\tr(A_zA_{z'})
  \le \frac1{L^2}\E_x\normF{C_q^{(i)}(x)}^2
  \begin{cases}
    1,&z=z',\\
    \theta,&z\ne z'.
  \end{cases}
\end{equation}
The new frequency Gram matrices thus have small total overlap in
expectation.  We can now use Lemma~\ref{lem:splitting} to assign them
to separate branches, charging the misplaced compressions to this
overlap.

\paragraph{Splitting the stack and updating the weights.}
For each $m$, apply Lemma~\ref{lem:splitting} to $(A_z)_{z\in\Lambda}$,
obtaining projections $(P_s)_{s\in\Lambda}$.  Define
\[
  H_{(q,s)}^{(i+1)}=P_sH'.
\]
Since $\sum_sP_s=I$ and $P_s^2=P_s=P_s^{\mathsf T}$, we have
$\sum_s(H_{(q,s)}^{(i+1)})^{\mathsf T}H_{(q,s)}^{(i+1)}
=H'^{\mathsf T}H'$.
Summing over $q$ proves \eqref{eq:partial-squares} at the new stage.

These projections may depend on $m_{\le i+1}$ and $x_{>i+1}$, but
they do not depend on any processed variable $y_{\le i+1}$.
Hence they commute with taking the Fourier coefficients in these
variables.  We retain all $R_{i+1}$ rows of every projected matrix,
even when the ranks of the projections vary with the
remaining inputs.  This gives a common coordinate system for every
branch at every value of those inputs.  Since all sets are finite,
the choices of projections for all histories and remaining inputs can
be made simultaneously, using no future averaging parameters.

The new distance satisfies
$\operatorname{dist}((q,s),(p,z))
=\operatorname{dist}(q,p)+\mathbf1_{\{s\ne z\}}$.
The Gram matrix with the new weights is therefore
\[
  C_{(q,s)}^{(i+1)}
  =\sum_{z\in\Lambda}\lambda^{\mathbf1_{\{s\ne z\}}}P_sA_zP_s.
\]
Cauchy--Schwarz for this sum of $L$ matrices, followed by
Lemma~\ref{lem:splitting}, gives
\begin{equation}\label{eq:one-step-potential}
\begin{split}
  \sum_s\normF{C_{(q,s)}^{(i+1)}}^2
  &\le L\left(
       \sum_z\normF{A_z}^2
       +\lambda^2C_L\sum_{z\ne z'}\tr(A_zA_{z'})
       \right).
\end{split}
\end{equation}
For the terms with $s=z$ we used
$\normF{P_zA_zP_z}\le\normF{A_z}$.
There are $L$ diagonal terms and $L(L-1)$ ordered off-diagonal terms
in \eqref{eq:overlap-summary}.  Taking expectation in
\eqref{eq:one-step-potential} yields
\[
  \E_m\sum_s\normF{C_{(q,s)}^{(i+1)}}^2
  \le\bigl(1+(L-1)\lambda^2C_L\theta\bigr)
          \E_x\normF{C_q^{(i)}(x)}^2
  \le2\E_x\normF{C_q^{(i)}(x)}^2.
\]

This estimate was proved for every fixing of $m_{\le i}$ and
$x_{>i+1}$.  Under the product probability law, the new $m$ is
independent of these variables.  Although earlier projections can
make $Z(x)$ an arbitrary function of the current input $x$, it is
independent of this new $m$, as required when applying
\eqref{eq:matrix-channel-variance}.  Thus we may sum over $q$ and
average over the fixed data to obtain $\Phi_{i+1}\le2\Phi_i$.
At $i=t$, this proves \eqref{eq:smoothing-squares} and
\eqref{eq:smoothing-potential}, with $R=R_t$.

\paragraph{From the Gram potential to Fourier mass.}
Choose $m$ with
$\sum_q\normF{C_q^m}^2\le2^tB^2$.
Each $U_q^m$ has $rL^t$ columns, so
$\rank(C_q^m)\le rL^t$, regardless of its row dimension.
For a positive semidefinite matrix $C$,
$\tr C\le\sqrt{\rank C}\,\normF{C}$.
Applying this inequality and then Cauchy--Schwarz over the $L^t$
branches gives
\[
\begin{split}
  \sum_{q,p}\lambda^{\operatorname{dist}(q,p)}
                     \normF{\widehat H_{q,p}^m}^2
  &=\sum_q\tr C_q^m\\
  &\le\sqrt{rL^t}\sum_q\normF{C_q^m}\\
  &\le\sqrt{rL^t}\sqrt{L^t}
            \left(\sum_q\normF{C_q^m}^2\right)^{1/2}\\
  &\le L^t\sqrt r\,B\,2^{t/2}.
\end{split}
\]
This is \eqref{eq:smoothing-coefficients}.
\end{proof}
